% ============================================================
%  Georg Brandes, Levned, Bind I
%  Kjøbenhavn: Gyldendal, 1905.
%  Pp. 129–132 (on Hans Brøchner)
%  TRANSLATION (English)
%  Translated by Hans Halvorson
% ============================================================
\documentclass[12pt,a4paper]{article}
\usepackage[utf8]{inputenc}
\usepackage[T1]{fontenc}
\usepackage[english]{babel}
\usepackage{microtype}
\usepackage{setspace}
\usepackage[margin=1.2in]{geometry}
\usepackage{fancyhdr}
\usepackage{hyperref}

\hypersetup{
  pdftitle={My Life, Vol. I — on Hans Brøchner},
  pdfauthor={Georg Brandes},
  hidelinks
}

\pagestyle{fancy}
\fancyhf{}
\rhead{Brandes, \emph{My Life} I (1905), pp.\ 129--132}
\cfoot{\thepage}

\setstretch{1.3}

% Original page-number marker
\newcommand{\pnum}[1]{\noindent\makebox[0pt][r]{\small\textit{[#1]}\quad}}

\title{\textbf{My Life} (\textit{Levned}), Volume I\\[6pt]
  \large Excerpt on Hans Brøchner (pp.\ 129--132)}
\author{Georg Brandes\\[4pt]
  \small Translated by Hans Halvorson}
\date{Copenhagen: Gyldendal, 1905}

\begin{document}
\maketitle
\thispagestyle{fancy}

\bigskip

\noindent For a long time I had pursued my non-juridical studies as
best I could without guidance from any teacher. But I had felt the want
of such guidance. And when a newly appointed lecturer at the University,
Professor H. Brøchner, offered guidance in the study of philosophy to
anyone who wished to present themselves at certain specified times, I
had felt strongly tempted to avail myself of the offer. I had had
reservations about it, for I was reluctant to surrender the least bit
of my precious freedom; I enjoyed my retired state, the mystery of my
modest student life; but on the other hand I would not be able to
venture into Plato and Aristotle without hints from an expert hand about
what and how.

I was therefore in great suspense. I had heard Professor Brøchner
lecture on psychology; but the delivery had been handled with such
anxious care, was so monotonous and alien-sounding, that it could not
but intimidate. Only the professor's distinguished beauty attracted,
especially his large sorrowful eyes with their lovely gaze, which
looked as though they had searched and wept.

\pnum{130}Now I resolved to venture up to Brøchner. But I had not
the courage to speak of it beforehand to my mother, lest I alarm myself
by naming it; so uneasy was I about the step. When the day came that I
had appointed for the visit --- it was the 2nd of September 1861 --- I
walked up and down before the house several times before I could bring
myself to mount the stairs; I tried endlessly to calculate in advance
what the professor would likely say and what I ought then best to reply.

The tall, handsome man, who resembled a Spanish knight, opened the
door himself and received kindly the young person who was soon to become
his closest student. To my astonishment he knew, the moment he heard my
name, which school I had come from and that I had recently become a
student. He strongly dissuaded a reading through of Plato and Aristotle
as far too strenuous --- said or implied that I ought to be spared the
laborious path he had himself traversed, and sketched a study plan for
modern philosophy and aesthetics. His bearing inspired confidence and
left the dominant impression that he wished to spare the beginner every
unnecessary effort. I felt, for all my young energies, nonetheless, as
I walked home, a little disappointed at not being to take hold of the
history of philosophy from the beginning.

The visit was soon repeated, and quickly there unfolded between teacher
and student the most intimate relationship --- from the elder's side a
fatherly benevolence such as the younger had not yet known, in which
the teacher, both instructing and practically guiding, kept his
attention continually fixed on everything that might promote the
pupil's welfare and secure his future; from the younger's side a
relationship of reverence and devotion, a deep gratitude for the care
of which he was the object.

I felt, to be sure, now and then before this teacher's great
\pnum{131}thoroughness and his ability to grapple with the most
difficult thoughts, a painful distrust of my own abilities and my own
intellectual power in comparison with his. I also not infrequently
felt uncomfortable when something like a crookedness seemed to enter
into our relationship --- when Brøchner, despite the doubts and
objections I raised, always took it as given that I must share his
pantheistic fundamental view, without noticing that I was still being
tossed about in doubts and groping for a foothold. But the intimate
relationship with the mature man had an attraction that the relationship
with uncertain or youthfully limited comrades necessarily had to lack:
he had a life's experience behind him, he looked from above down upon a
young person's sympathies and antipathies.

For me, for instance, Ploug's \textit{Fædrelandet} still stood at
that time as Denmark's intelligent organ, while Bille's
\textit{Dagbladet} was distasteful to me, especially on account of the
superficiality and the dismissive tone that characterized the
newspaper's literary criticism. Brøchner, who with mixed benevolence
and without making much distinction looked down upon both of these
newspapers of the educated public, found in his young pupil's
indignation at the trivial yet useful newspaper only a mark of his
disposition. But merely the way in which Brøchner one day remarked with
a smile --- ``I suppose you don't read \textit{Dagbladet} as a matter
of principle'' --- this made me see in a flash the comical quality in
my indignation at a Hjerrild's articles. My horizon was still narrow
enough that conditions in Copenhagen seemed, in and of themselves,
significant. I myself considered my horizon broad. One day, when I took
stock of my inner accounts, I wrote, as a nineteen-year-old, naïvely:
``My good qualities, those that will constitute my personality if I
amount to something, are a powerful and glowing enthusiasm, a capable
authority in the service of truth, a wide \pnum{132}horizon and a
philosophically educated thinking. These must reconcile one to my lack
of humor and ease.''

In essential agreement with Hans Brøchner I first felt myself several
years after the acquaintance was formed. But by then a study of Ludwig
Feuerbach's writings had swept me away; for Feuerbach was the first
thinker in whom I found the origin of the idea of God in the human mind
satisfactorily explained. In Feuerbach I finally encountered an
exposition without circumlocution and without the usual heavy formulas
of German philosophy, a victorious clarity that had a highly salutary
effect on my own manner of thought and gave me security. If for several
years I had felt myself considerably to the right of my friend and
teacher, there now came a time when I felt myself considerably to the
left of him with his mysterious life in the eternal realm of spirit, of
which he felt himself a part.

\end{document}
